\subsection{The maximal fixed base of the chosen arithmetic parameter}

\begin{statement}{D.49}{Proved}{A simultaneous fixed-base lift, the entire congruence action, and a remaining mixed kernel}
Keep the actual arithmetic element \(h\), its normalized parameter
\(f_h\), and the complete D.48 action. All subgroups and commutators
below are closed. Let \(C_h\) be the closed normal closure of
\(f_h\) in \(\widehat B\), and define
\[
 \begin{gathered}
 C_h=\overline{\langle f_h\rangle}^{\,\widehat B},\qquad
 \mathcal H_h=\widehat B/C_h,\\
 \mathfrak R_{h,p}=\ker(C_h\to C_h^{(p)}),\qquad
 \mathfrak R_h=\bigcap_p\mathfrak R_{h,p}.
 \end{gathered}
 \tag{D.49.1}
\]
Then \(C_h\subset M^+\). The marked quotient \(\mathcal H_h\)
is exactly the maximal profinite quotient on which this fixed
\(\alpha_h\) acts pointwise as identity. Put
\[
 \mathcal P_{h,p}=\widehat B/\mathfrak R_{h,p},\qquad
 \mathcal W_h=\widehat B/\mathfrak R_h.
 \tag{D.49.2}
\]
There are actual marked identifications
\[
 \begin{gathered}
 \mathcal P_{h,p}
   =\varprojlim_{K}\widehat B/\ker(K\to K^{(p)}),\qquad
 \mathcal W_h=\prod_{\mathcal H_h}\mathcal P_{h,p},\\
 \ker(\mathcal W_h\to\mathcal H_h)=\prod_pC_h^{(p)},\qquad
 \ker(\mathcal W_h\to\mathcal P_{h,p})
                       =\prod_{q\ne p}C_h^{(q)}.
 \end{gathered}
 \tag{D.49.3}
\]
Here \(K\) runs through all open normal subgroups of
\(\widehat B\) with \(C_h\subset K\subset L\), where
\(L=[\widehat N,\widehat B]\) is D.48's actual subgroup.
Equivalently, these are all finite \(h\)-fixed marked bases
refining \(\mathcal E\). This is one compatible inverse
system, with both source-kernel inclusions proved below.
The common marked quotient of distinct prime factors is
exactly \(\mathcal H_h\).

On this source there is a continuous peripheral automorphism
\[
 \begin{gathered}
 \alpha_h^*=(\alpha_{h,h,2},(1_{\mathcal P_{h,p}})_{p\ne2}),
 \\
 f_h^*=(f_{h,h,2},(1)_{p\ne2})
                  \in\operatorname{im}(C_h\to\mathcal W_h).
 \end{gathered}
 \tag{D.49.4}
\]
The double subscript in \(\alpha_{h,h,2}\) means the action
of the original arithmetic \(h\) on the new source
\(\mathcal P_{h,2}\). The multiplier is one, and one pair
of normalized conjugators is constructed on the entire
product. The action and parameter project exactly to D.48.
In particular the prescribed CM pair, all Fermat homology
actions, and the full metabelian nonarithmetical outer class
are retained.

The common quotient and projection kernels retain the
non-solvable contribution of the new base. Define
\[
 \begin{aligned}
 T_h&=M^+/C_h,\\
 \mathfrak r(T_h)&=\ker\bigl(T_h\to\prod_pT_h^{(p)}\bigr),\\
 \mathcal D_h&=\mathcal H_h/\mathfrak r(T_h)
                 =\widehat B/(C_hR^+),\\
 D_{h,p}&=\ker((M^+)^{(p)}\to T_h^{(p)}),\\
 E_{h,p}&=\ker(C_h^{(p)}\twoheadrightarrow D_{h,p}).
 \end{aligned}
 \tag{D.49.5}
\]
Then \(\mathcal D_h\) is the exact maximal common marked
quotient of \(\mathcal H_h\) and \(\mathcal V^+\), and
\[
 \begin{gathered}
 1\longrightarrow\prod_pE_{h,p}\longrightarrow\mathcal W_h
   \longrightarrow
       \mathcal H_h\times_{\mathcal D_h}\mathcal V^+
   \longrightarrow1,\\
 1\longrightarrow\prod_pE_{h,p}
   \longrightarrow\ker(\mathcal W_h\to\mathcal V^+)
   \longrightarrow\mathfrak r(T_h)\longrightarrow1,\\
 \mathfrak R_h\subset R^+.
 \end{gathered}
 \tag{D.49.6}
\]
These are exact sequences with their original source maps;
no splitting is asserted. The group \(\mathfrak r(T_h)\)
has an actual \(A_5\) quotient under the marked
\(\operatorname{PSL}_2(\mathbb F_5)\) specialization.
Thus the new source is strictly larger than D.48's source.

In fact the full congruence quotient, with its single
global sign, factors through the new identity component:
\[
 \widehat B\longrightarrow\mathcal P_{h,3}
        \twoheadrightarrow
 \mathcal Q_{\rm cong}
       :=\operatorname{SL}_2(\widehat{\mathbb Z})/\{\pm I\}.
 \tag{D.49.7}
\]
Consequently \(\alpha_h^*\) acts pointwise as identity on
\(\mathcal Q_{\rm cong}\). For the original arithmetic
\(h\), the exact statements used are
\[
 \kappa_{2,\ell}(h)=0\quad(\ell\ne3),\qquad
 \kappa_{2,3}(h)\in3\mathbb Z_3.
 \tag{D.49.8}
\]
Full three-adic vanishing for that original \(h\) is not
asserted. The factorization (D.49.7), with its pro-three
kernel calculation below, proves the entire congruence
assertion without imposing that additional vanishing.

There is also an explicit remaining finite mixed source.
Set
\[
 \begin{gathered}
 G=G_3=\widehat B/(\widehat N'\widehat N^3),
 \qquad K_3=\widehat N'\widehat N^3,\qquad n=|G|=576\cdot3^{97},\\
 \mathcal J_{3,2}=\widehat B/(K_3'K_3^2),\qquad
 V_{3,2}=K_3/(K_3'K_3^2)\simeq\mathbb F_2^{\,1+n/6}.
 \end{gathered}
 \tag{D.49.9}
\]
The actual element \(a=f_h\bmod G\) is nonzero of order
three in \(\widehat N/(\widehat N'\widehat N^3)\).
Let \(A=\langle a^G\rangle\) be its actual normal closure
in that elementary abelian three-group, and define the
submodule using the actual deck action:
\[
       W=[V_{3,2},A]
          =\sum_{b\in A}(b-1)V_{3,2}.
 \tag{D.49.10}
\]
The exact common marked quotient of
\(\mathcal W_h,\mathcal J_{3,2}\) is
\(\mathcal J_{3,2}/W\), and
\[
 \begin{gathered}
 \operatorname{im}(\mathfrak R_h\to\mathcal J_{3,2})=W,\qquad
 \dim_{\mathbb F_2}W={n\over6}\left(1-{1\over|A|}\right),\\
 (a-1)V_{3,2}\subset W,\qquad
 \dim_{\mathbb F_2}(a-1)V_{3,2}={n\over9}=64\cdot3^{97}>0.
 \end{gathered}
 \tag{D.49.11}
\]
The exact formula uses the actual finite subgroup \(A\).
Its dimension has not been numerically evaluated; the
displayed nonzero cyclic subspace and its dimension are
already determined. Thus \(\mathfrak R_h\) is nontrivial,
and the remaining full-source lift is not replaced by a
cofinality claim. Genuine full profinite GT0 realization
of the CM pair and of the retained metabelian action
remains undecided.

Finally the attempted repetition of D.48's central step
is completely decided:
\[
             [L,\widehat B]=L,\qquad
       \widehat B/[L,\widehat B]\xrightarrow{\ \sim\ }
                                      \widehat B/L.
 \tag{D.49.12}
\]
The exact integral augmentation map on the 193-cycle
lattice has Smith form \((1^{193})\); its Hermite image
is the full lattice. This calculation and the maps above
continue Part D(e), section 51.16, folio 734 verso,
scan 174, and the two-torus/Fermat follow-ups in
section 51.15, folios 725--727, scans 156--160.
\end{statement}
\begin{proof}
\emph{The exact maximal fixed base.}
For the original arithmetic element, the normalized
identities are
\(\alpha_h(s)=s f_h\),
\(\alpha_h(r)=f_h^{-1}r\), and
\(\alpha_h(t)=t\).
A marked quotient map agrees with its composite with
\(\alpha_h\) exactly when it kills \(f_h\): necessity
follows from the first equality, and sufficiency follows
from both generator equalities. This proves the maximal
property of \(\mathcal H_h\) and pointwise invariance.
It also proves \(\alpha_h(C_h)=C_h\), including the
reverse inclusion by composing with \(\alpha_h^{-1}\).
D.48 proves that \(h\) fixes \(\mathcal H^+_\infty\),
so \(C_h\subset M^+\subset\widehat F'\).

For the first inverse limit in (D.49.3), restriction
\(C_h\to K^{(p)}\) proves one residual-kernel inclusion.
For the other, start with the kernel of any finite
continuous \(p\)-quotient of \(C_h\). Uniform continuity
on the compact group \(C_h\) gives a finite conjugation
orbit; its intersection is a \(\widehat B\)-normal open
subgroup \(C_0\) with finite \(p\)-group quotient.
Choose open normal \(O\) in \(\widehat B\) with
\(O\cap C_h=C_0\). The open normal subgroups \(K/C_h\)
with \(K\subset L\) are cofinal neighbourhoods of one
in \(\mathcal H_h\). Some such \(K\) lies in \(OC_h\);
its image modulo \(O\) is a \(p\)-group. Its residual
kernel consequently lies in \(O\). The intersection of
all those kernels already lies in \(C_h\), and hence
lies in every \(C_0\). This proves the reverse inclusion
and the actual inverse-limit equality.

For distinct primes, the maximal pro-prime images of
\(C_h\) have no common nontrivial finite quotient.
Finite Goursat followed by compactness gives their full
product as the image of \(C_h\). This proves every
product, common-quotient and projection-kernel assertion
in (D.49.3), before choosing any automorphism.

\emph{One coherent peripheral automorphism and its projection.}
Because the actual \(h\) fixes the entire base
\(\mathcal H_h\), its two-primary action and the
odd-primary identity actions define an automorphism
of the actual fibre product. Let \(A_h,\beta_h\)
be its original \(r,s\)-conjugators. Their images in
\(\mathcal H_h\) centralize \(r,s\), respectively.
On the two-primary factor, multiply \(A_h\) on the
right by an \(r\)-centralizing lift of its inverse
base coordinate. Such a lift exists by the coprime
order-three fixed-coset argument through the pro-two
kernel. Take the identity \(r\)-conjugator on every
odd factor. All these base coordinates are now one.
For \(s\), retain \(\beta_h\) at two, and on each odd
factor choose an \(s\)-centralizing lift of the same
whole base coordinate. Coprime order-two lifting
through the pro-\(p\) kernel supplies these elements.
The finite complement argument and compactness used
in D.46 apply to these entire profinite base elements,
so the choices are compatible at every finite level.

The resulting tuples lie in \(\mathcal W_h\). Their
parameter is exactly \(f_h^*\) in (D.49.4), with an
actual preimage in \(C_h\subset\widehat F'\).
They fix \(t\), give the two hexagons with multiplier
one, and generate because the product action is an
automorphism. This proves the continuous peripheral
construction, not only its finite linear equations.

Restriction \(C_h\to(M^+)^{(p)}\) gives
\(\mathfrak R_{h,p}\subset R_p^+\) for every prime.
Thus \(\mathfrak R_h\subset R^+\), with actual maps
\(\mathcal P_{h,p}\to\mathcal P_p^+\) and
\(\mathcal W_h\to\mathcal V^+\). Each prime component
of (D.49.4) maps to the corresponding D.48 component;
the action and parameter therefore project exactly.
The images of the two new conjugators differ on the
right from any already retained D.48 pair by elements
of the respective centralizers of \(r,s\).
The two actions agree on all of \(\mathcal V^+\),
and D.25's \(G_2/G_3\) rigidity gives trivial \(Q\)-image
for these differences. This proves retention of the
normalized CM coordinates through their actual maps.
It also proves all the claimed homology and metabelian
comparisons by projection, without replacing the
retained arithmetic representative.

\emph{The common quotient with D.48, including its full kernel.}
For a surjection of compact groups \(P\to T\), the
image of the pronilpotent residual of \(P\) is the
pronilpotent residual of \(T\). Indeed the quotient
by that image is a quotient of the maximal
pronilpotent quotient of \(P\), so is pronilpotent;
in the other direction the maximal pronilpotent
quotient of \(T\) is a pronilpotent quotient of
\(P\). Compactness makes the image closed.
Apply this to \(M^+\to T_h\). The image of
\(R^+\) in \(\mathcal H_h\) is exactly
\(\mathfrak r(T_h)\), giving both descriptions of
\(\mathcal D_h\) in (D.49.5).

The original two quotient kernels \(C_h\) and
\(R^+\) have product \(C_hR^+\). Closed Goursat
therefore identifies their full joint source image
with \(\mathcal H_h\times_{\mathcal D_h}\mathcal V^+\).
For each prime the image of \(C_h\) in
\((M^+)^{(p)}\) is exactly \(D_{h,p}\).
To prove this, divide by that closed normal image:
the resulting group has the universal property of
the maximal pro-\(p\) quotient of \(M^+/C_h\).
The image of \(C_h\) in the product of all these
prime groups is the full product \(\prod_pD_{h,p}\),
again by coprime finite Goursat and compactness.
Its new product image is \(\prod_pC_h^{(p)}\).
The kernel of the joint map is consequently
\(\prod_pE_{h,p}\), proving the first exact sequence
in (D.49.6). Projecting the fibre product to
\(\mathcal V^+\) gives kernel
\(\mathfrak r(T_h)\), which proves the second
sequence. In particular that kernel is not silently
replaced by only the product of the prime kernels.

\emph{The full congruence factor and its global sign.}
Let
\(\mathcal Q^{(3')}
 =(\prod_{\ell\ne3}\operatorname{SL}_2(\mathbb Z_\ell))
 /\{\pm I\}\), with one global sign on that product.
The original \(s,r\) matrices give the actual
surjection \(\widehat B\to\mathcal Q_{\rm cong}\).
At finite level surjectivity follows from elementary
row reduction over each local ring \(\mathbb Z/\ell^k\)
and Chinese remainders: upper and lower unipotents
generate, and are the actual powers of the marked
braid matrices. Passing to the finite inverse system
gives the stated map with its single global sign.

For \(\tau\in I_3\) and \(\ell\ne3\), every
\(\ell\)-power root of unity and every \(\ell\)-power
root of the unit two lies in the maximal unramified
extension at three. Reduction to the algebraic residue
closure followed by Hensel's lemma proves this, since
the derivative is prime to three. Hence
\(\chi_\ell(\tau)=1\) and
\(\kappa_{2,\ell}(\tau)=0\). The actual commutator
formula of D.46 gives \(\kappa_{2,\ell}(h)=0\)
for every \(\ell\ne3\). At three its exact value is
\[
 (1-\chi_3(\tau))\kappa_{2,3}(g)
       +(\chi_3(g)-1)\kappa_{2,3}(\tau)\in3\mathbb Z_3,
 \tag{D.49.13}
\]
using the conditions actually imposed in D.46.
This proves precisely (D.49.8).

D.18's normalized matrix formula now shows that
the original \(h\) fixes \(\mathcal Q^{(3')}\)
pointwise. It also fixes \(Q_{12}=Q\).
Since \(\alpha_h(s)=s f_h\), both quotient maps
kill \(f_h\), and hence their joint kernel contains
the full congruence image of \(C_h\).
That joint kernel is exactly
\[
 \begin{aligned}
 \ker(\mathcal Q_{\rm cong}\to
                  \mathcal Q^{(3')}\times Q_{12})
   &\simeq\ker\bigl(\operatorname{SL}_2(\mathbb Z_3)
                           \to\operatorname{SL}_2(\mathbb F_3)\bigr),\\
 b&\longmapsto[((I)_{\ell\ne3},b)].
 \end{aligned}
 \tag{D.49.14}
\]
For both inclusions, an element in the away-from-three
kernel has a representative with every away component
equal to \(I\), after one global sign. The retained
modulo-four coordinate is then \(I\), so identity
in \(Q_{12}\) forces the modulo-three coordinate
also to be \(I\), with positive sign. Conversely
every displayed principal-congruence element is
in the joint kernel. The map is injective because
that away identity representative cannot be changed
by a nontrivial global sign. The successive three-adic
congruence layers are additive three-groups, so the
kernel is pro-three.

Thus the actual full congruence image of \(C_h\)
is pro-three. Its source map kills
\(\mathfrak R_{h,3}\), proving (D.49.7).
The identity action on \(\mathcal P_{h,3}\)
therefore proves the full congruence assertion.
The certificate independently checks all 576
\(Q_{12}\) states: projection to the independent
local projective groups has image of order 288
and kernel \(\{I,5I\}\). Thus the relative-sign
factor detected in (D.49.14) is retained, including
the finite check that prevents losing it.

The \(\operatorname{PSL}_2(\mathbb F_5)\) specialization
is a quotient of \(\mathcal H_h\). The program
enumerates its 60 matrices and its five Klein-four
subgroups. Conjugation gives a faithful even
permutation action on those five subgroups, hence
the exact marked \(A_5\) map. The actual permutations
of \(s,r\), in cycle notation on \(\{0,1,2,3,4\}\), are
\((1\,3)(2\,4)\) and \((0\,4\,3)\).
Its conjugacy-class sizes \(1,12,12,15,20\)
also verify simplicity directly: the only unions
containing identity whose orders divide 60 have
orders one and 60. The image of \(M^+\) in this
simple quotient is all of \(A_5\), because its
quotient is a quotient of the prosolvable
\(\mathcal H^+_\infty\). As \(C_h\) is killed,
\(T_h\) maps onto \(A_5\), and its pronilpotent
residual does too. This proves the asserted strict
enlargement and non-solvable kernel contribution.

\emph{An exact nonzero remaining mixed kernel.}
D.47 proves that the original arithmetic
\(r\)-conjugator has \(Q\)-image \(5I\).
D.25 proves
\(q(C_{G_3}(r)\cap\operatorname{im}\widehat F)=1\).
Therefore the original \(h\) cannot act as identity
on \(G_3\). If its actual parameter had trivial
image there, its two normalized generator formulas
would make that action identity. Hence
\(a=f_h\bmod G_3\) is nonzero. It lies in the
elementary abelian kernel of \(G_3\to Q\), so
has exact order three. This establishes the actual
source element used in (D.49.10).

The kernel \(K_3\) is a torsion-free free subgroup
of index \(n\) in \(B\), with rank \(1+n/6\).
Its actual modulo-two graph homology gives
\(\mathcal J_{3,2}\) and \(V_{3,2}\).
Let \(H\) be the image of \(C_h\) in this finite
quotient. It is the normal closure of the actual
image of \(f_h\), and projects onto \(A\).
Commuting its conjugates with the whole kernel
\(V_{3,2}\) shows \(W\subset H\).
The quotient \(H/W\) is a central extension of
the elementary abelian three-group \(A\) by an
elementary abelian two-group. Coprime complements
split this extension, so it is abelian.

Conversely averaging over \(A\) in characteristic
two gives \([W,A]=W\): the augmentation image
has zero coinvariant quotient. A three-Sylow
complement to \(H\cap V_{3,2}\) in \(H\)
projects isomorphically onto \(A\). In every
nilpotent quotient its three-primary image commutes
with the two-primary image of \(W\), forcing all
of \(W\) to vanish. Thus the exact nilpotent
residual of \(H\) is \(W\). The residual-image
argument already proved above, now applied to
\(C_h\to H\), gives the first equality in
(D.49.11). The product of the two actual source
kernels then gives the common quotient
\(\mathcal J_{3,2}/W\).

For its dimension, every nonidentity element of
\(A\) has trivial image in \(Q\), whereas every
nonidentity vertex stabilizer in the Bass--Serre
graph has a nontrivial order-two or order-three
image in \(Q\). Thus \(A\) acts freely on all
edges and vertices of the graph for \(K_3\).
Since its order is odd, the category of
\(\mathbb F_2[A]\)-modules is semisimple by
averaging. The actual graph chain sequence
therefore gives the module decomposition
\[
 V_{3,2}\simeq\mathbb F_2
          \oplus\mathbb F_2[A]^{\,n/(6|A|)}.
 \tag{D.49.15}
\]
For clarity, this follows by taking the edge
permutation module, subtracting the two vertex
permutation modules, and adding the constant
\(H_0\) module in the split graph sequence:
their free-module ranks are respectively
\(n/|A|\), \(n/(2|A|)\), and \(n/(3|A|)\).
Semisimplicity makes the resulting equality of
constituents an actual module isomorphism.
The augmentation part has dimension
\((|A|-1)n/(6|A|)\), proving the formula for
\(W\). Apply the same actual graph argument to
the specific cyclic subgroup \(\langle a\rangle\)
of order three. Its moving subspace
\((a-1)V_{3,2}\) has dimension \(n/9\), and lies
in \(W\). This proves the explicit nonzero
remaining source-kernel image without evaluating
the undetermined dimension of the whole \(Q\)-span.

\emph{The fully tested central route and its stabilization.}
The program \texttt{d\_cm\_fixed\_parameter.py}
reads the actual D.48 marked quotient and rebuilds
its integral 193-cycle lattice. It reconstructs
the two left generator permutations, verifies that
they commute with every right marked generator,
and verifies their action on every full integral
edge chain. The augmentation matrix has shape
\(193\times386\). Its Smith diagonal is
\((1^{193})\), and independently its Hermite
image is \(I_{193}\). Thus
\(L^{\rm ab}/((1-s)L^{\rm ab}+(1-r)L^{\rm ab})=0\).
Because \([L,\widehat B]\) contains \(L'\),
completion gives \([L,\widehat B]=L\).
This proves the actual quotient isomorphism
(D.49.12); a further central repetition cannot
produce the new mixed lift constructed above.
\end{proof}
